% Generated draft: requires mathematical and source review.
\documentclass[11pt]{article}
\usepackage[margin=25mm]{geometry}
\usepackage{fontspec}
\setmainfont{Latin Modern Roman}
\usepackage{amsmath,amssymb,mathtools,unicode-math}
\setmathfont{Latin Modern Math}
\usepackage{enumitem,needspace,xurl,hyperref}
\setlength{\parindent}{0pt}
\setlength{\parskip}{5pt}
\setlength{\emergencystretch}{3em}
\newcommand{\N}{\mathbb{N}}
\newcommand{\Z}{\mathbb{Z}}
\newcommand{\Q}{\mathbb{Q}}
\newcommand{\R}{\mathbb{R}}
\newcommand{\C}{\mathbb{C}}
\newcommand{\F}{\mathbb{F}}
\newcommand{\A}{\mathbb{A}}
\newcommand{\PP}{\mathbb{P}}
\newcommand{\E}{\mathbb{E}}
\newcommand{\eps}{\varepsilon}
\newcommand{\dd}{\,\mathrm{d}}
\begin{document}
\section*{UnsolvedMath ID: 30000986}
\textbf{Source problem code:} catalogue entry
\subsection{Definitions and mathematical statement}
Let $P$ be a finite $p$-group with the source's generalized Cayley map onto a Lie algebra $\mathfrak u$. Let $\operatorname{Irr}(P)$ be its irreducible complex characters, and let $\mathcal O$ be the coadjoint-orbit set for the required compatible action on the dual of $\mathfrak u$. The extract does not define the generalized Cayley map, action, coefficient field, or dual convention. For an additive-character dual, the coadjoint action would be $(g\cdot\ell)(x)=\ell(\operatorname{Ad}(g^{-1})x)$. The original asks for conditions giving, for every integer $k\geq0$, \[\{\chi\in\operatorname{Irr}(P):\chi(1)=p^k\}=\{O\in\mathcal O:|O|=p^{2k}\}.\] Literally this compares characters with orbits, which are different kinds of objects. The natural fake-degree interpretation is equality of their counts, \[\bigl|\{\chi:\chi(1)=p^k\}\bigr|=\bigl|\{O:|O|=p^{2k}\}\bigr|\qquad(k\geq0),\] or existence of a degree-compatible bijection, not set identity without a specified correspondence. This interpretation is explicitly flagged rather than silently changing the displayed source. The original says the unrestricted conjecture is false for $p=2$ and known for algebra groups over odd primes. The later supplied 2026-08-21 assessment reports counterexamples in every characteristic and positive results for important odd-prime algebra groups; the scope of the original broad positive claim therefore needs clarification.\par Separately let $U_n(q)$ be the upper unitriangular group over $\F_q$, and let $h_{U_n(q)}$ be its number of conjugacy classes. For each fixed $n$, the related polynomiality question is whether there exists $t_n\in\Z[T]$ such that \[h_{U_n(q)}=t_n(q)\qquad\text{for every prime power }q.\] The polynomial may depend on $n$; no single polynomial independent of $n$ is asserted. The original describes an equivalent commuting-variety question. In the usual unitriangular model let $\mathfrak u_n(\F_q)$ be the strictly upper triangular matrices and \[\mathcal C_n(\F_q)=\{(X,Y)\in\mathfrak u_n(\F_q)^2:XY=YX\}.\] The normalization relating the counts is $|\mathcal C_n(\F_q)|=|U_n(q)|h_{U_n(q)}=q^{n(n-1)/2}h_{U_n(q)}$. Thus the source's point-count formulation must retain this group-order factor. The exact fixed-$n$ and characteristic range is not explicitly stated in the extract. The supplied triage labels the combined record partially solved and the general conditions question open.
\subsection{Short English statement}
Ask when irreducible-character degrees are counted by square roots of coadjoint-orbit sizes, and whether unitriangular class numbers are integer-coefficient polynomials in the field size for each fixed matrix size. The literal source set equality and conflicting odd-characteristic status claims require clarification.
\subsection{Sources}
\begin{itemize}\item Local UnsolvedMath registry, permanent ID 30000986, source problem catalogue entry, \emph{Fake Degrees and Class Numbers of Unitriangular Groups}, proposed year 2008; no individual proposer supplied.\item Supplied citation: \emph{Profinite and Asymptotic Group Theory} (2009), Oberwolfach Reports, DOI 10.4171/owr/2008/28, pp.\ 1537--1588: \url{https://doi.org/10.4171/owr/2008/28}. Supplied literature label: original OWR report.\item Supplied literature label: counterexamples in all characteristics, \url{https://arxiv.org/abs/1502.03242}.\end{itemize}
\end{document}
